Source : Eurostat, comptes annuels des administrations publiques (gov\_10a\_main, gov\_10a\_taxag, gov\_10a\_exp), PIB (nama\_10\_gdp) ; Insee, dette trimestrielle de Maastricht par sous-secteur (contributions consolidées, 4e trimestre). Calculs de l'auteur (dette/attribution.py).\par Lecture : fin 2025, les 702 Md\euro{} de dette de 1995, capitalisés au taux apparent des APU (6,5 \% en 1996, 1,3 \% en 2020), représentent 1 928 Md\euro{} ; les déficits primaires de 1996 à 2025 et leurs intérêts, 1 355 Md\euro{}. C'est une décomposition comptable, pas causale : sans excédents primaires, les intérêts de la dette héritée sont payés par de nouveaux emprunts.
